\chapter{Внутренний перенос энтропии и предел четырёхсостояний часов}

Предыдущая проверка показала, что проекторная среда может распасться после
глубокого переохлаждения, но не объяснил рост обратной температуры.
Теперь проверяется минимальная замкнутая возможность: может ли семейный
триплет охладиться, передав энергию и энтропию уже существующему
четырёхсостояний носителю внутренних часов.

Вопрос имеет два независимых уровня. На кинематическом уровне требуется
унитарная орбита полного состояния, на которой редуцированное состояние
триплета становится одноосным. На динамическом уровне тот же переход
должен сохранять полную энергию, быть автономным и задавать устойчивую
стрелу внутреннего времени.

\section{Точный энтропийный бюджет кристаллизации}

Изотропное семейное состояние имеет физическую энтропию
\begin{equation}
 S_{\rm iso}=\log3=1.0986122887\ldots.
\end{equation}
В точке сосуществования предыдущей проверки
\begin{equation}
 R_*=\operatorname{diag}
 (0.9121665963\ldots,0.0439167019\ldots,0.0439167019\ldots)
\end{equation}
и
\begin{equation}
 S_*=0.3583777189\ldots.
\end{equation}
Следовательно, для упорядочения необходимо вывести из триплета
\begin{equation}
 \boxed{\Delta S_{\rm out}=S_{\rm iso}-S_*
 =0.7402345698\ldots.}
 \label{eq:v6-ordering-entropy-export}
\end{equation}

Эффективная энергия фаз меняется на
\begin{equation}
 \Delta E_{\rm rel}=E_{\rm iso}-E_*
 =0.4798400112\ldots,
 \label{eq:v6-ordering-energy-release}
\end{equation}
и в точке сосуществования выполнено машинно проверенное тождество
\begin{equation}
 \beta_c\Delta E_{\rm rel}=\Delta S_{\rm out}.
 \label{eq:v6-coexistence-entropy-energy-identity}
\end{equation}
Это не новый закон, а равенство свободных энергий двух фаз.

\section{Четыре состояния имеют достаточную ёмкость}

Чистая $d$-мерная подсистема способна увеличить свою энтропию не более чем
на $\log d$. Для четырёхсостояний носителя
\begin{equation}
 S_{C,\max}=\log4=1.3862943611\ldots,
\end{equation}
что почти вдвое больше \eqref{eq:v6-ordering-entropy-export}. Требуемая
доля ёмкости равна
\begin{equation}
 \frac{\Delta S_{\rm out}}{\log4}=0.5339663714\ldots.
\end{equation}
Минимальная чистая принимающая система должна иметь размерность
\begin{equation}
 d_{\min}=\left\lceil e^{\Delta S_{\rm out}}\right\rceil=3.
\end{equation}
Таким образом, проектные четыре состояния проходят первый нетривиальный
бюджетный тест без добавления новых степеней свободы.

\begin{proposition}[Энтропийная достаточность четырёх состояний]
Четырёхсостояний носитель способен однократно принять всю энтропию,
теряемую одним семейным триплетом при переходе из $I_3/3$ в фазу
$R_*$. Размерностного запрета на внутреннее охлаждение нет.
\end{proposition}

\section{Явная замкнутая унитарная орбита}

Пусть часы начинаются в чистом состоянии $|0\rangle_C$, а триплет --- в
$I_3/3$:
\begin{equation}
 \rho_{\rm in}=\frac{I_3}{3}\otimes|0\rangle\langle0|.
 \label{eq:v6-entropy-transfer-initial-state}
\end{equation}
Для спектра $r=(r_0,r_1,r_2)=\Spec R_*$ определим равновесный ансамбль
\begin{equation}
 |\phi_k\rangle
 =\sum_{j=0}^2\sqrt{r_j}\,e^{2\pi i kj/3}|j\rangle,
 \qquad k=0,1,2.
\end{equation}
Он удовлетворяет
\begin{equation}
 \frac13\sum_{k=0}^2|\phi_k\rangle\langle\phi_k|=R_*.
\end{equation}
Рассмотрим полное состояние
\begin{equation}
 \rho_{\rm out}=\frac13\sum_{k=0}^2
 |\phi_k\rangle\langle\phi_k|\otimes|k\rangle\langle k|_C.
 \label{eq:v6-entropy-transfer-target-state}
\end{equation}
Его редукции равны
\begin{equation}
 \rho_{S,\rm out}=R_*,
 \qquad
 \rho_{C,\rm out}=\frac13\sum_{k=0}^2|k\rangle\langle k|.
\end{equation}
При этом $\rho_{\rm in}$ и $\rho_{\rm out}$ имеют один глобальный спектр
$(1/3,1/3,1/3,0,\ldots,0)$. Поэтому существует полный унитарий
$W\in U(12)$ такой, что
\begin{equation}
 \rho_{\rm out}=W\rho_{\rm in}W^*.
\end{equation}

Это конструктивно доказывает: замкнутая унитарная система способна сделать
редуцированный триплет более упорядоченным, не уменьшая полной энтропии.
Энтропия переносится в часы и взаимные корреляции.

Для начального произведения точный баланс имеет вид
\begin{equation}
 \Delta S_C-\Delta S_{\rm out}=I(S:C)_{\rm out}\geq0.
 \label{eq:v6-subsystem-entropy-balance}
\end{equation}
В явной конструкции
\begin{equation}
 S_{C,\rm out}=\log3,
 \qquad
 I(S:C)_{\rm out}=S_*=0.3583777189\ldots.
\end{equation}

Каноническая типичность действительно позволяет малой подсистеме выглядеть
тепловой внутри чистого или стационарного большого целого; см.
\path{arXiv:quant-ph/0511225} и \path{doi:10.1038/nphys444}.
Общие условия приближения подсистем к равновесию исследованы в
\path{arXiv:0812.2385} и \path{doi:10.1103/PhysRevE.79.061103}.
Наш результат слабее этих теорем: он доказывает существование одной
унитарной орбиты, а не типичность или релаксацию.

\section{Цена: принимающая подсистема перестаёт быть чистыми часами}

До переноса носитель был чист. После переноса он имеет ранг три и энтропию
$\log3$. В выбранной конструкции он хранит запись о трёх компонентах
ансамбля, но теряет чистую фазовую когерентность. Значит, один и тот же
четырёхсостояний объект не может без обратной реакции одновременно быть
идеальными часами и бесплатным стоком энтропии.

Литература по автономным квантовым часам подчёркивает, что часы являются
неравновесным ресурсом и их точность связана с энтропийным производством;
см. \path{arXiv:1609.06704}. Конечные автономные часы также испытывают
обратную реакцию при управлении машиной; см. \path{arXiv:1607.04591}.
Это точно соответствует найденному проектному ограничению.

\begin{nogocriterion}[Запрет бесплатных часов-резервуара]
Нельзя считать четырёхтактный носитель одновременно неизменным эталоном
времени и стоком энтропии. Если он принимает энтропию кристаллизации, его
редуцированное состояние и качество часов обязаны измениться.
\end{nogocriterion}

\section{Почему условие четвёртого порядка ещё не задаёт холодильник}

Проект строго построил унитарный шаг $U$ с $U^4=I$ и четырьмя характерными
проекторами. Но из одного $U$ не следует единственный гамильтониан часов.
Любой логарифм содержит выбор ветвей,
\begin{equation}
 H_C=\frac1{\tau}
 V\operatorname{diag}(\theta_j+2\pi n_j)V^*,
 \qquad n_j\in\mathbb Z,
 \label{eq:v6-clock-logarithm-ambiguity}
\end{equation}
где $e^{-i\theta_j}$ --- собственные значения $U$.
Не фиксированы ни $\tau$, ни целые $n_j$, ни абсолютная энергия.

Следовательно, пока нельзя проверить условие
\begin{equation}
 [W,H_S+H_C+H_{\rm int}]=0
\end{equation}
или сопоставить выделившуюся энергию
\eqref{eq:v6-ordering-energy-release} с возбуждением часов. Явный унитарий
доказывает доступность состояния, но может требовать внешней работы.

Самосодержащие квантовые холодильники существуют, однако используют
несколько энергетически резонансных подсистем и тепловых ресурсов; см.
\path{doi:10.1103/PhysRevLett.105.130401}. Это показывает математическую
возможность автономного охлаждения, но не выводит нужную машину из одного
четырёхтактного сдвига проекта.

\section{Корреляции и стрела времени}

Глобальная унитарность не запрещает локальное охлаждение, но и не выбирает
его направление. Начальные корреляции способны даже обратить обычный
тепловой поток; см. \path{doi:10.1103/PhysRevE.77.021110} и
\path{doi:10.1103/PhysRevE.81.061130}. Поэтому специальную начальную
запутанность нельзя принимать как бесплатное объяснение стрелы: она сама
является ресурсом и граничным условием.

Кроме того, конечномерная автономная унитарная система рекуррентна. Она
может однократно перевести энтропию из триплета в часы, но не создаёт строго
монотонную $\beta(\tau)$ на неограниченном интервале. Для эффективной
необратимости нужен большой носитель, многократные внутренние ячейки,
термодинамический предел или явно обоснованное крупнозернистое описание.

\begin{theorem}[Кинематическое внутреннее охлаждение]
Проектный четырёхсостояний носитель имеет достаточную энтропийную ёмкость,
и существует замкнутая унитарная орбита, переводящая изотропный семейный
триплет в состояние сосуществования $R_*$. Однако существующий
четырёхтактный сдвиг не фиксирует энергию, взаимодействие и автономный
закон перехода; перенос также портит сами часы. Поэтому внутреннее
охлаждение открыто динамически, но больше не запрещено кинематически.
\end{theorem}

\begin{protocol}
\begin{enumerate}
 \item энтропия изотропного триплета: \textbf{$\log3$};
 \item энтропия упорядоченной фазы при сосуществовании:
 \textbf{$0.3583777189\ldots$};
 \item необходимый экспорт энтропии:
 \textbf{$0.7402345698\ldots$};
 \item высвобождаемая эффективная энергия:
 \textbf{$0.4798400112\ldots$};
 \item ёмкость четырёх чистых состояний: \textbf{$\log4$};
 \item минимальная размерность чистого стока: \textbf{три};
 \item четырёхсостояний носитель достаточен: \textbf{да};
 \item явная глобальная унитарная орбита: \textbf{построена};
 \item сохранение чистых часов после переноса: \textbf{нет};
 \item гамильтониан из одного $U^4=I$: \textbf{не определяется};
 \item энергосохраняющее автономное взаимодействие: \textbf{не выведено};
 \item монотонная внутренняя $\beta(\tau)$: \textbf{не выведена};
 \item следующая проверка:
 \textbf{управляемый часами энергосохраняющее переключение};
 \item рождение материи:
 \textbf{энтропийно допустимо, динамически не замкнуто}.
\end{enumerate}
\end{protocol}

Машинный сертификат сохранён в
\path{s2t_v6_internal_entropy_transfer_cooling_gate_results.json}.